Following \cite[Definition 5.20]{io}, we define a \emph{slice} of $\tqdn$ to be a full subquiver $S$ of $\tqdn$ such that:
\begin{enumerate}
\item Any $\nu_{d}$ orbit in $\tqdn$ contains precisely one vertex which belongs to $S$.
\item $S$ is convex, i.e., for any path $P$ in $\tqdn$ connecting two vertices in $S$, all vertices appearing in $P$ belong to $S$.
\end{enumerate}
Slices are shown in blue in Figure~\ref{fig:q13} and Figure~\ref{fig:q23}.

Given a slice $S$ of $\tqdn$, one can find a cut $C_{S}$ of $\qdn$ such that $Q_{C_{S}}^{(d, n)}$ is isomorphic to $S$ with the arrows $F_{i}$ of $Q_{C_{S}}^{(d, n)}$ corresponding to arrows $A \to A + E_{i}$ in $S$ \cite[Theorem 5.24]{io}.
